 \section{Green-hyperbolic operators}\label{sec:GHOs}

{\bf Preliminaries.}
 We begin by recalling the general setting of Green-hyperbolic operators~\cite{Baer:2015}. Let~$M$ be a smooth finite-dimensional manifold, allowing the possibility that $M$ has finitely many connected components with possibly different dimensions, and let $g$ be a smooth Lorentzian metric on $M$ of signature $+--\cdots$. With these structures, $M$ is automatically Hausdorff and paracompact~\cite{Geroch_spinorI:1968}.
 We assume that $(M,g)$ is time-orientable and that a time-orientation has been chosen. To minimise notation, we denote the Lorentzian spacetime formed by the manifold, metric and time orientation with the single symbol $M$. The volume measure induced by the metric will be denoted $\mu$. On other points of notation,
 the symbol $\subset$ will always allow for the possibility of equality, while $\NN_0$ and $\NN$ denote the natural numbers with or without zero, respectively.

 As usual, the causal future/past of a point $x\in M$ is denoted $J^\pm(x)$ and comprises all points (including $x$) that may be reached from $x$ along smooth future/past-directed curves. (Throughout this paper, we tacitly order alternatives labelled by $\pm$ or $\mp$ so that the alternative labelled by the upper symbol comes first.)
 If $S\subset M$ then one writes $J^\pm(S)=\cup_{x\in S} J^\pm(x)$, and $J(S)=J^+(S)\cup J^-(S)$. The spacetime is globally hyperbolic if it contains no causal curves and $J^+(K)\cap J^-(K)$ is compact for all compact sets $K\subset M$~\cite{Bernal:2006xf}. Globally hyperbolic spacetimes can be foliated into smooth spacelike Cauchy surfaces~\cite{Bernal:2004gm}; it is also the case that $J^\pm(K)$ are closed whenever $K$ is compact. We adopt the following terminology: $S\subset M$ is
 \emph{spacelike compact} if $S$ is closed and $S\subset J(K)$ for some compact $K$; $S$ is \emph{future/past-compact} if $S\cap J^{\pm}(x)$ is compact for all $x\in M$; $S$ is \emph{strictly future/past-compact} if
 $S\subset J^{\mp}(K)$ for some compact $K$.

 If $B$ is a vector bundle with finite-dimensional fibres, then $\Gamma^\infty(B)$ will denote the corresponding space of smooth sections and $\Gamma_{0/\pc/\fc/\spc/\sfc/\mathrm{sc}}^\infty(B)$ will be the spaces of smooth sections with com\-pact/past-compact/future-compact/strictly past-compact/strictly future-compact/space\-li\-ke-compact support. The space of smooth sections with support contained in some a closed subset $A\subset M$ is denoted $\Gamma^\infty_A(M)$. Further details on the topologies of these spaces are summarised in Appendix~\ref{appx:topspaces}. We will only consider bundles with finite-dimensional fibres that are (without any real loss) vector spaces over $\CC$. The bilinear (not sesquilinear!) pairing between sections of dual bundle $B^*$ and sections of $B$ is denoted with angle brackets $\langle \cdot,\cdot\rangle \colon \Gamma^\infty(B^*)\times \Gamma^\infty(B)\to C^\infty(M)$.

{\bf Green-hyperbolicity.} Suppose that $B_1$ and $B_2$ are bundles over $M$ and that $P\colon \Gamma^\infty(B_1)\to \Gamma^\infty(B_2)$ is a linear partial differential operator. Then there is a formal dual $\adj{P}\colon \Gamma^\infty(B_2^*)\to \Gamma^\infty(B_1^*)$ given by
 \begin{equation*}
 	\int_M \mu\, \big\langle \adj{P}f, \phi\big\rangle = \int_M\mu\, \langle f, P\phi\rangle
 \end{equation*}
 for all $\phi\in \Gamma^\infty(B_1)$ and $f\in \Gamma^\infty(B_2^*)$ whose supports intersect compactly; $\adj{P}$ is also a linear partial differential operator.
 If they exist, linear maps $E^\pm\colon \GoinX{B_2}\to \Gamma^\infty(B_1)$ obeying
 \begin{itemize}\itemsep=0pt
 	\item[(G1)] $E^\pm P f=f$ for all $f\in\GoinX{B_1}$,
 	\item[(G2)] $P E^\pm f = f$ for all $f\in\GoinX{B_2}$,
 	\item[(G3)] $\supp E^\pm f \subset J^\pm (\supp f)$ for all $f\in\GoinX{B_2}$
 \end{itemize}
 are called advanced ($-$) and retarded ($+$) Green operators for $P$.\footnote{B\"ar reverses the usage of `advanced' and `retarded'; we adopt the more standard convention.}
 \begin{Definition}
 	$P$ is said to be \emph{Green-hyperbolic} if both $P$ and $\adj{P}$ admit advanced and retarded Green operators.
 \end{Definition}

 It is a remarkable fact that this purely algebraic definition -- only requiring linearity of $E^\pm$ and with no presumption of uniqueness or continuity -- has the following consequence, which is a summary of results in~\cite[Sections 3 and 4]{Baer:2015} and deserves to be called the \emph{first main theorem of Green-hyperbolicity}.
 \begin{Theorem}\label{thm:GHbasics}
 	Let $P\colon\Gamma^\infty(B_1)\to \Gamma^\infty(B_2)$ be a Green-hyperbolic operator and consider:
 	\begin{enumerate}\renewcommand{\theenumi}{(\roman{enumi})}\itemsep=0pt
 		\item[$(i)$] the restriction $P\colon\Gamma^\infty_{\spc/\sfc}(B_1)\to\Gamma^\infty_{\spc/\sfc}(B_2)$,
 		\item[$(ii)$] the restriction $P\colon\Gamma^\infty_{\pc/\fc}(B_1)\to\Gamma^\infty_{\pc/\fc}(B_2)$, or
 		\item[$(iii)$] the extension $P\colon\DD'_{\pc/fc}(B_1)\to\DD'_{\pc/fc}(B_2)$ $($see below$)$.
 	\end{enumerate}
 	Then, in each case, $P$ has continuous inverses, denoted
 	$\tilde{E}^\pm$, $\overline{E}^\pm$ and $\widehat{E}^\pm$ in cases
 	$(i)$, $(ii)$, $(iii)$ respectively, and which are successive extensions of $E^\pm$. Each of these inverses has the support property $(G3)$, replacing
 	$\Gamma_0^\infty(B_2)$ by the appropriate domain. In particular,
 	$E^\pm$ are uniquely determined by $P$ and are continuous.
 \end{Theorem}	
 In $(iii)$ the space of distributional sections $\DD'(B)$ of a bundle $B$ over $M$ is defined
 as the topological dual of $\Gamma_0^\infty(B^*\otimes\Omega)$, where $\Omega$ is the bundle of weight-$1$ densities over $M$. As sections of $B$ and $B^*\otimes\Omega$ can be paired to give a density, there is an obvious embedding of $\Gamma^\infty(B)$ in $\DD'(B)$, with $\phi\in\Gamma^\infty(B)$ corresponding to the distribution
 $f\mapsto \int_M \langle \phi,f\rangle$ acting on $f\in \Gamma_0^\infty(B^*\otimes\Omega)$.
 The map $P$ in $(iii)$ is the restriction of the dual map
 $\big(\mu(\adj{P})\mu^{-1}\big)'$ to elements of $\DD'(B_1)$ with past- or future-compact support.
 (Recall that the formal dual is defined relative to the specific density~$\mu$.)
 In~\cite{Baer:2015}, B\"ar defines distributional sections of $B$ using
 the topological dual of $\Gamma_0^\infty(B^*)$, which would be the space of distribution densities in our terminology. The metric density
 provides an isomorphism between the spaces and operators that he considers and those that we do. As we have in mind potential applications where more than one metric might be in use, we have elected not to make a fixed identification between distributions and distribution densities.

